\section{Rational Euler sources and repeated quadratics}\label{sec:quadratic}
One way to build a small linear form is to cancel many initial terms of
a generating series, leaving a rapidly decaying tail. A rational
correction is attractive because integration by parts can turn its
contribution into a rational boundary value. The Euler operator below
multiplies $t^k$ by $k$; applying it four times matches the fourth powers
in the beta series. We first reduce such a construction to a polynomial
integral, then test one natural repeated-factor family. The obstruction
comes from fourth powers of primes that survive in the reduced form.

Let $E=t\,\mathrm d/\mathrm dt$, $g(t)=t/(1+t^2)$, and
\[
 J(R)=\frac16\int_0^1(-\log t)^3R(t)\frac{\dd t}{t}.
\]
If $H$ is rational, regular on $[0,1]$, and $H(t)=O(t)$ at zero,
put $h(z)=H(e^{-z})$. Then $E=-\mathrm d/\mathrm dz$, and
$J(E^4H)=\frac16\int_0^\infty z^3h^{(4)}(z)\dd z$.
Four integrations by parts leave $h(0)=H(1)$; the other boundary
terms at zero contain $z^3$, $z^2$ or $z$, and those at infinity
vanish by exponential decay. Since $J(g)=\theta$, this gives
\begin{equation}\label{eq:euler-exact}
 J(E^4H)=H(1),\qquad J(Ag+E^4H)=A\theta+H(1).
\end{equation}
This isolates the target exactly, but imposes no favorable arithmetic
bound on the coefficient $A$.

Put $F(y)=\sum_{k\geq0}(-1)^ky^k/(2k+1)^4$. Let
$Q(y)=\sum_{j=0}^d q_jy^j$ have no zero on $[0,1]$, and let $C_m$
be the Taylor polynomial of $QF$ through degree $m-1$. The choice
$H(t)=-AtC_m(t^2)/Q(t^2)$ cancels the first $m$ odd coefficients
of $Ag+E^4H$: indeed $E^4(tF(t^2))=g(t)$, and
$t(F(t^2)-C_m(t^2)/Q(t^2))$ vanishes to order at least $2m+1$.
To see what remains at $t=1$, define
\[
 W_{Q,m}(y)=\sum_{j=m}^d q_j+
                 \sum_{j=0}^{\min(m-1,d)}q_j(-y)^{m-j}.
\]
For a term $q_jy^j$ with $j<m$, the factor $F$ contributes the
omitted tail $I((-t)^{m-j})$ at $y=1$; for $j\geq m$, the
whole factor $F$ contributes $I(1)$. Weighting by $q_j$ and summing gives
$Q(1)\theta-C_m(1)=I(W_{Q,m})$, hence the exact formal-pair identity
\begin{equation}\label{eq:reversal}
 \theta-\frac{C_m(1)}{Q(1)}=\frac{I(W_{Q,m})}{Q(1)}.
\end{equation}
This reduction includes $m<d$ and introduces no period-independence
assumption. For $Q=(v+uy)^n$ and $m=2n$,
$W_{Q,2n}=(-y)^n(u-vy)^n$. The rational Euler source therefore has
the same primitive pair, up to sign, as the following scalar source.

\begin{theorem}[Repeated-quadratic sources; S03]\label{thm:quadratic}
For even $n$, let $u_n,v_n$ be positive integers with
$\log(u_n+v_n)=o(n)$. Set
$P_n(y)=y^n(u_n-v_ny)^n$ and let $L_n>0$ be the value of its
primitive formal pair under~\eqref{eq:formal-pair}. Then
\[
 \liminf_{\substack{n\to\infty\\n\ {
m even}}}
 \frac{\log L_n}{n}
 \geq\frac83-\log(5+4\sqrt2)>0.
\]
\end{theorem}
\begin{proof}
The proof has two parts: primes in $(4n/3,2n)$ force a denominator
even after gcd removal, and a minimax estimate limits how small the
normalized positive integral can be.
Write $A=P_n(-1)=(u+v)^n$ and $I(P_n)=A\theta+N/D$ in lowest
terms. The first $n$ coefficients of $(P_n-A)/(1+y)$ are
$(-1)^{j+1}A$. For a prime
\[
 \frac{4n}{3}<p<2n,\qquad p\nmid(u+v),
\]
only the denominator $p$ among $1,3,\ldots,4n-1$ is divisible by $p$.
Its index $(p-1)/2<n$ has a $p$-unit coefficient. Consequently
$v_p(N/D)=-4$, exactly. If $g=\gcd(DA,N)$, then $g\mid A$, so
\begin{equation}\label{eq:quadratic-primitive}
 L_n=\frac Dg I(P_n)
 \geq\prod_{\substack{4n/3<p<2n\\p\nmid(u+v)}}p^4\,
                  \frac{I(P_n)}A.
\end{equation}
This controls the actual primitive value, including all common content.
The gcd step uses $\gcd(D,N)=1$, which implies $g\mid A$ and
$D/g\geq D/A$. The selected prime factors are unavoidable factors
of $D$, whereas $A$ accounts for the most that gcd removal could save.

Let $r=v/(u+v)$ and $f_r(y)=y(1-r-ry)$. With
$\dd\mu=\dd\nu/(1+y)$,
$I(P_n)/A=\int|f_r(y)|^n\dd\mu(y)$. An elementary minimax calculation gives
\[
 M_*:=\min_{0\leq r\leq1}\max_{0\leq y\leq1}|f_r(y)|
       =\frac1{5+4\sqrt2}.
\]
Indeed for $r\geq1/3$ the interior positive maximum is
$(1-r)^2/(4r)$; for $r\geq1/2$ the negative endpoint has magnitude
$2r-1$. They are equal at $r=(1+2\sqrt2)/7$.
For $r<1/3$ the maximum is at least $1/3$.

Uniformity for moving $r$ follows from $|f_r'|\leq2$. In an interval
of length $\ell=M_*/(4n)$ adjacent to a maximizer,
$|f_r|\geq M(r)(1-1/(2n))$. On its central half,
$1-y\geq\ell/4$ and $\dd\mu\geq(1-y)^3\dd y/192$.
Since $(1-1/(2n))^n\geq1/2$,
\[
 \frac{I(P_n)}A\geq\frac{M_*^{n+4}}{3\cdot2^{22}n^4}.
\]
The prime number theorem in~\eqref{eq:quadratic-primitive} now gives
\[
 \log L_n\geq
 \left(\frac83-\log(5+4\sqrt2)\right)n
 -4\sum_{\substack{4n/3<p<2n\\p\mid u+v}}\log p+o(n).
\]
The exceptional sum is at most $\log(u+v)=o(n)$.
\end{proof}

It suffices more generally that the displayed exceptional prime sum
be $o(n)$. The theorem does not cover odd $n$, different osculation
lengths, general denominator polynomials, or sums of these forms.
The source-level statement is weaker: if $H=P/Q$ has integer
coefficients and $p\nmid Q(0)$, cancellation of the first $m$ odd
coefficients forces
$v_p(A)\geq4\max_{k<m}v_p(2k+1)$. This follows coefficientwise from
$A(-1)^k+(2k+1)^4[t^{2k+1}]H=0$. It does not by itself control
primitive endpoint content; that was the purpose of the prime argument.
